%------------------------------------------------------------------------------------
%  SECTION 1 : ANNOTATIONS ET SUIVI DES MODIFICATIONS
%------------------------------------------------------------------------------------
\section{Annotations et Notes de Suivi}
\label{sec:todos}

Durant la phase de rédaction, il est essentiel de pouvoir annoter le document sans polluer le texte définitif.
L'option de classe \texttt{todoNote} active l'affichage des notes marginales et de la table des révisions.

\subsection{Notes dans la marge et surlignage}
\label{subsec:notes_marge}

La commande personnalisée \texttt{\textbackslash todoH} permet d'associer directement un surlignage coloré dans le corps du texte à une note en marge :

Dans cette phrase, nous identifions un point à relire avec attention \todoH{MarieStyle}{Emerald!30}{Vérifier les données statistiques}{les données cliniques présentées ici}.
De plus, une seconde remarque peut être formulée pour l'auteur \todoH{GuillaumeStyle}{Apricot!40}{Ajouter la référence manquante}{concernant l'algorithme d'imagerie médicale}.

\subsection{Gestion des figures provisoires}
\label{subsec:missing_figures}

Lorsqu'une figure n'est pas encore finalisée, la commande \texttt{\textbackslash missingfigure} réserve l'espace dans la mise en page sans bloquer la compilation :

\missingfigure[figwidth=\linewidth, figheight=4.5cm, figcolor=RoyalBlue!20]{Schéma de synthèse du protocole d'étude à insérer}

